\documentclass[11pt,a4paper]{article}
\usepackage[T1]{fontenc}
\usepackage{lmodern,microtype,amsmath,amssymb,amsthm}
\usepackage[margin=25mm]{geometry}
\usepackage{booktabs,array,xurl}
\usepackage[hidelinks,unicode]{hyperref}
\hypersetup{pdftitle={Erdos 708: An Attempt at the Joint Inequality},pdfauthor={Research working note}}
\setlength{\emergencystretch}{3em}
\newtheorem{theorem}{Theorem}[section]
\newtheorem{lemma}[theorem]{Lemma}
\newtheorem{proposition}[theorem]{Proposition}
\newtheorem{corollary}[theorem]{Corollary}
\theoremstyle{remark}\newtheorem{remark}[theorem]{Remark}
\newcommand{\pos}[1]{(#1)_+}
\newcommand{\LP}{L_{\rm raw}}
\newcommand{\atLP}{L_{\rm at}}
\newcommand{\one}{\mathbf1}
\title{Erd\H{o}s 708: An Attempt at the Joint Inequality\\
\large Exact loss under deletion of weight blocks,\newline sparse-obstruction reductions, and the remaining gap}
\author{Research working note}
\date{23 November 2026}
\begin{document}
\maketitle
\begin{abstract}
The requested inequality \(J_{\rm weak}\), and hence the unrestricted
asymptotic bound \(g(n)\le(2+o(1))n\), are \textbf{not proved in this note}.
The main proved result is an exact joint loss inequality for an optimal
boxed covering dual: for disjoint coordinate blocks, the sum of the
objective losses caused by deleting each block separately is at most the
original optimum. It yields a genuine gap between small and substantial
weights at a controlled objective cost, without assuming such a gap.
We also prove a sparse-optimal mean bound and a weighted reduction to a
bounded number of substantial occurrences per retained source. These are
necessary reductions of a hypothetical general obstruction, rather than
new source-class covering theorems. A verified arithmetic counterexample
rules out one same-threshold comparison shortcut; it does not refute
\(J_{\rm weak}\). We state the precise surviving gluing obligation and
keep a stronger prime-power-count LP separate from the original problem.
\end{abstract}
\tableofcontents

\section{The requested statement and the status of the argument}

Let \(A\subseteq[2,M]\cap\mathbb N\) consist of \(m\) distinct integers,
and let \(I=\{x+1,\ldots,x+M\}\), \(x\ge0\). Put
\[
 D_p(A)=\sum_{a\in A}v_p(a),\qquad
 X_w(n)=\sum_p w_pv_p(n),\qquad H_I(w)=\sum_{b\in I}\pos{X_w(b)-1}.
\]
The source primes form a finite set. The boxed valuation LP is
\begin{equation}
 \LP(A,I)=\min\left\{\sum_{b\in I}z_b:
 0\le z_b\le1,\quad\sum_bz_bv_p(b)\ge D_p(A)\ \text{for every }p\right\}.
 \label{eq:primal}
\end{equation}
Its dual objective is
\[
 \Phi_A(w)=\sum_{a\in A}X_w(a)-H_I(w).
\]
An optimum may be chosen with \(0\le w_p\le1\). All valuations here are
raw; target exponents are not truncated at source exponents or at \(M\).

For the original requested minimal class, fix \(\delta>0\), set
\(c=2+3\delta\), and take a least-cardinality integer failure at the same
\(M,I\):
\[
 \tau(A,I)>cm,\qquad \tau(S,I)\le c|S|\quad(S\subsetneq A).
\]
The preceding report \cite{previous} established the following
normalization. If \(w\) is optimal, \(r\) is the number of source primes,
and \(G=\{a:X_w(a)>2+2\delta\}\), then
\begin{gather}
 r=o_\delta(m),\quad |A\setminus G|=o_\delta(m),\quad
 W:=\sum_{a\in A\setminus G}X_w(a)=o_\delta(m),\label{eq:oldnorm}\\
 \Phi_G(w)>c|G|-r.\label{eq:oldlower}
\end{gather}
The union \(U\) of the retained prime-deletion banks is saturated in one
full optimal primal solution, and
\[
 \Xi=\sum_p w_p\pos{V_p(U)-D_p(G)}\le W=o_\delta(m),\qquad
 V_p(U)=\sum_{b\in U}v_p(b).
\]
The exact requested inequality is
\begin{equation}
 \boxed{\quad
 \sum_{a\in G}\bigl(X_w(a)-2-\delta\bigr)
 \le H_I(w)+C_\delta\Xi+\delta|G|+o_\delta(m).
 \quad}\tag{$J_{\rm weak}$}\label{eq:jweak}
\end{equation}
The error must be uniform in \(M,A,x,w\) on the full stated class.

If \eqref{eq:jweak} held, then
\(\Phi_G(w)\le(2+2\delta)|G|+o_\delta(m)\), contradicting
\eqref{eq:oldlower}. Equivalently, restore \(W\) once and round the raw
LP at an extra cost \(r=o(m)\). This gives
\(\tau(A,I)\le(2+2\delta)m+o(m)<cm\).
The known ambient gate \(M\ge16m^2\Rightarrow\tau(A,I)\le2m\)
ensures that the minimal sizes tend to infinity in any unbounded
counterexample sequence \cite{wang,previous}. Thus the missing estimate
would settle the unrestricted asymptotic question. Exact \(2n\) would
still require a separate argument.

This note does not establish \eqref{eq:jweak}. In particular, none of the
necessary restrictions proved below may be substituted for its missing
joint hinge bound.

\section{An exact joint bound for deleting coordinate blocks}

This section applies to \emph{every} optimal boxed valuation LP, without
a minimality, sparsity, core-size, or prime-gap assumption.
Choose paired optimal solutions \(z,w\), with \(0\le w_p\le1\), and write
\[
 X=X_w,\quad F=\LP(A,I),\quad H=\{b\in I:X(b)>1\},\quad
 T=\{b\in I:X(b)=1\}.
\]
Complementary slackness gives
\begin{equation}
 z_b=1\ (b\in H),\qquad z_b=0\ (X(b)<1),\qquad
 \sum_bz_bv_p(b)=D_p(A)\ (w_p>0).
 \label{eq:cs}
\end{equation}
Consequently
\begin{equation}
 F=|H|+\sum_{b\in T}z_b.\label{eq:Fboundary}
\end{equation}
The threshold-one contribution is essential and is retained below.

\begin{theorem}[Joint block-deletion loss]
Let \(B_1,\ldots,B_N\) be pairwise disjoint sets of prime coordinates;
they need not exhaust the support. Let \(w^{(j)}\) be obtained by setting
the weights on \(B_j\) to zero. Put
\[
 \Delta_j(n)=\sum_{p\in B_j}w_pv_p(n),\qquad
 \ell_j=F-\Phi_A(w^{(j)}).
\]
Then
\begin{equation}
 \ell_j=\sum_{b\in H}\pos{\Delta_j(b)-(X(b)-1)}
       +\sum_{b\in T}z_b\Delta_j(b),\label{eq:blockexact}
\end{equation}
and
\begin{equation}
 \boxed{\quad 0\le\sum_{j=1}^N\ell_j\le F.\quad}
 \label{eq:blocksum}
\end{equation}
In particular, one block can be deleted at a loss of at most \(F/N\).
\end{theorem}

\begin{proof}
Every nonzero contribution to \(\Delta_j\) belongs to an active tight
row, so
\[
 \sum_{a\in A}\Delta_j(a)=\sum_{b\in I}z_b\Delta_j(b).
\]
Insert this in the difference of the two dual objectives. At a hot
target its contribution is
\[
 \Delta_j-\bigl[(X-1)-(X-\Delta_j-1)_+\bigr]
 =\pos{\Delta_j-(X-1)}.
\]
At a threshold-one target the contribution is \(z_b\Delta_j(b)\).
At a cold target both hinges vanish and \(z_b=0\). This proves
\eqref{eq:blockexact} and the nonnegativity of each loss.

For nonnegative numbers \(y_j\) and \(t\ge0\),
\(\sum_j(y_j-t)_+\le(\sum_jy_j-t)_+\). At each hot target,
disjointness gives \(\sum_j\Delta_j(b)\le X(b)\), and hence
\[
 \sum_j\pos{\Delta_j(b)-(X(b)-1)}\le1.
\]
At a threshold-one target the sum is at most \(z_bX(b)=z_b\).
Summing and using \eqref{eq:Fboundary} proves \eqref{eq:blocksum}.
\end{proof}

This proof spends each target's capacity once. Its bound contains no
factor for the number of source patterns, the number of primes, or
the total weighted source load. The statement concerns deleting each
block separately from the same optimum; it does not justify adding
these loss bounds for the simultaneous deletion of many blocks.

\subsection{A weight gap derived at a controlled cost}

Fix \(R>1\) and an integer \(N\ge1\), and take
\[
 B_j=\{p:R^{-j-1}<w_p\le R^{-j}\},\qquad 0\le j<N.
\]
For one of these bands, write \(b=R^{-j}\), \(a=b/R\), and let \(w'\)
be the deleted vector. Then
\begin{equation}
 \Phi_A(w')\ge F-F/N,\qquad
 w'_p\in[0,a]\cup(b,1].\label{eq:weightgap}
\end{equation}
This is a gap in the weights, not an assumed gap between prime sizes.
For fixed \(N,R\), the selected threshold belongs to a finite fixed
list, so uniform statements valid at each fixed threshold can be
applied to the selection.

The sparse reduction below gives at most
\(K(b)\le4(\lfloor1/b\rfloor+1)\) substantial occurrences per retained
source. Thus
\begin{equation}
 a(K(b)+1)\le9/R.\label{eq:gapsmall}
\end{equation}
This avoids the invalid parameter choice of taking \(K=O(1/\eta)\)
and then assuming that \(\eta K\) becomes small as \(\eta\to0\).
It does not bound the total number of substantial prime coordinates.

More generally the bands may be chosen recursively as
\((b_{j+1},b_j]\), with \(1=b_0>b_1>\cdots>b_N>0\).
After \(b_j\) is fixed, \(b_{j+1}\) may be made as small as any
specified function of \(K(b_j)\). The same loss bound selects one
of the finitely many bands. This resolves parameter requirements on
the light weight itself, but does not bound a later multiplier on
the objective loss \(F/N\).

\section{Minimality and the exact costs of renormalization}

For some arguments it is useful to choose a separate raw-minimal
failure at fixed \(M,I\), with budget \(c>2\):
\[
 F=\LP(A,I)=cm+e>cm,\qquad
 \LP(S,I)\le c|S|\quad(S\subsetneq A).
\]
Such a reduction is justified in the preceding report. It is distinct
from the original integer-minimal class in Section 1.
At an optimal dual, weak duality on \(A\setminus\{a\}\) gives
\[
 X_w(a)\ge c+e>c,\qquad
 0<e\le\log_2M-c,
\]
and the quadratic ambient gate gives \(M<16m^2\). In particular
\(F=cm+o(m)\) and \(|H|=O_c(m)\).

For any nonnegative test vector \(v\), proper-subset minimality gives
the exact global shortfall bound
\begin{equation}
 \sum_{a\in A}\pos{c-X_v(a)}\le\pos{cm-\Phi_A(v)}.
 \label{eq:shortfall}
\end{equation}
Indeed, take \(S=\{a:X_v(a)>c\}\). If \(S=A\) the assertion is
immediate. Otherwise \(\Phi_S(v)\le c|S|\), with the same full target
hinge, and subtraction gives \eqref{eq:shortfall}.

Apply this to a band deletion with loss \(\ell\le F/N\). For
\(0<\gamma<c\), let
\(B_{\rm low}=\{a:X_{w'}(a)\le c-\gamma\}\). Then
\begin{equation}
 |B_{\rm low}|\le\ell/\gamma,\qquad
 W_{\rm low}:=\sum_{a\in B_{\rm low}}X_{w'}(a)
 \le(c-\gamma)\ell/\gamma.\label{eq:lowdiscard}
\end{equation}
If a further exceptional set has original weighted demand \(o(m)\),
delete it too and call the remaining set \(G'\), of size \(g'\). Then
\begin{equation}
 g'\ge m-F/(N\gamma)-o(m),\qquad
 \Phi_{G'}(w')\ge cg'-F/N-o(m).\label{eq:aftergap}
\end{equation}
The second inequality retains the cardinality credit from deleting
low-load sources; it is sharper than merely subtracting all their
weighted demand.

Although \(w'\) need not be optimal, \(H'\subseteq H\), and every
remaining active row is still exactly tight in the original primal
solution. Thus, for every \(U'\subseteq H'\),
\begin{equation}
 V_p(U')\le D_p(A)\quad(w'_p>0),\qquad
 \sum_pw'_p\pos{V_p(U')-D_p(G')}
 \le \frac{(c-\gamma)F}{N\gamma}+o(m).\label{eq:aftergaprow}
\end{equation}
This error is generally a small \emph{linear} error, not \(o(m)\) at
fixed \(N\). If a subsequent estimate multiplies it by a constant
depending on \(K(b)\), that dependence must be checked: increasing
\(N\) can also decrease \(b\) and increase \(K(b)\). No such constant
may be absorbed by an unsupported choice of ``sufficiently large
\(N\).''

\section{Uniform constraints from a sparse hot set}

Write \(\rho=m/M\) and
\(\mu=\sum_p w_p/(p-1)\). In this section sparsity means a sequence of
actual optimal obstructions with \(\rho\to0\); it is not asserted that
all obstructions have this property.

\subsection{Cold initial segment}

If \(X(d)>1\), every interval multiple of \(d\) lies in \(H\), so
\[
 \lfloor M/d\rfloor\le h:=|H|,\qquad d>M/(h+1).
\]
Consequently \(X(d)\le1\) for every integer \(d\le M/(h+1)\).
For fixed \(\eta>0\), set \(s=\lfloor1/\eta\rfloor+1\). If
\(w_p\ge\eta\), then
\begin{equation}
 p^s>M/(h+1).\label{eq:coldprime}
\end{equation}
Thus substantial primes tend to infinity in the sparse regime, since
\(h=O_c(m)\). This conclusion uses raw additivity; it must not be
transferred to capped per-prime profiles in Section 7.

\subsection{An exponential moment and the weighted mean}

Let \(t=\log(3/2)\). A nonnegative divisor expansion gives
\begin{equation}
 \frac1M\sum_{n\le M}e^{tX(n)}
 \le\prod_p\frac{p-1}{p-e^{tw_p}}\le e^\mu.
 \label{eq:expmoment}
\end{equation}
For completeness, the local divisor coefficients are
\(g(p^j)=e^{tw_pj}-e^{tw_p(j-1)}\ge0\). Counting multiples bounds the
mean by \(\sum_dg(d)/d\), the displayed Euler product. Since
\(e^{tw}-1\le w/2\) and
\(p-e^{tw}\ge(p-1)/2\), the logarithm of each factor is at most
\(w/(p-1)\).
Distinctness and Jensen's inequality now imply
\begin{equation}
 \sum_{a\in A}X(a)\le\frac m t\{\log(M/m)+\mu\}.
 \label{eq:actualjensen}
\end{equation}

All source primes are at most \(M\). Prime-power multiple counts give
\begin{equation}
 0\le\mu-\frac1M\sum_{n\le M}X(n)\le e_M,
 \qquad e_M=\frac{3\pi(M)+\sqrt M\log_2M}{M}=O(1/\log M).
 \label{eq:prefixmean}
\end{equation}
For one prime with \(J_p=\lfloor\log_pM\rfloor\), the rounding and
infinite-tail errors total at most \((J_p+2)/M\). Summing and counting
the prime powers proves the explicit bound. This argument does not
truncate any target exponent.

Finally,
\(H_I(w)\ge\sum_I X-M\ge\sum_{n\le M}X-M\). Combining this with
\eqref{eq:actualjensen}--\eqref{eq:prefixmean} yields, when \(\rho<t\),
\begin{equation}
 (1-\rho/t)\mu\le1-F/M+(\rho/t)\log(1/\rho)+e_M.
 \label{eq:meanbound}
\end{equation}
Hence along every sparse raw-minimal sequence,
\begin{equation}
 \boxed{\quad\mu\le1+O(\rho\log(1/\rho))+O(1/\log M)=1+o(1).\quad}
 \label{eq:sparsemean}
\end{equation}
The analogous integer-minimal bound adds only its known rounding
error divided by \(M\).

\subsection{A joint bound for substantial primes in a lower range}

Fix \(\eta\in(0,1]\), let \(s\eta>1\) with integer \(s\ge2\), and put
\[
 P_\eta(z)=\{p\le z:w_p\ge\eta\},\quad k=|P_\eta(z)|,
 \quad\mu_0=\sum_{p\in P_\eta(z)}1/p,
\]
\[
 E_s(k)=\binom{k}{s}+(s+1)\binom{k}{s+1}.
\]
The following finite inequality holds for every shift:
\begin{equation}
 \min(\mu_0,1/2)\le
 s\left(\frac{h+1}{M}\right)^{1/s}
 +\left(\frac{2s!\{h+E_s(k)\}}M\right)^{1/s}.
 \label{eq:bonf}
\end{equation}
Here is the proof, retaining the arithmetic error. Let
\(a_0=\max_{p\in P_\eta(z)}1/p<((h+1)/M)^{1/s}\), and choose a
subset \(Q\) of reciprocal mass \(\mu'\le1/2\) with
\(\mu'\ge\min(\mu_0,1/2)-a_0\). For every nonnegative integer \(N\),
\[
 \one_{N\ge s}\ge\binom Ns-(s+1)\binom N{s+1}.
\]
Apply this to the number of distinct \(Q\)-primes dividing a target.
Using ordinary multiple counts and the elementary symmetric sums
\(e_j((1/p)_{p\in Q})\) gives
\[
 h+E_s(k)\ge M(e_s-(s+1)e_{s+1})\ge\tfrac12Me_s.
\]
The second step uses \((s+1)e_{s+1}\le\mu'e_s\). Ordered products give
\[
 s!e_s\ge\pos{\mu'-(s-1)a_0}^{s}
 \ge\pos{\min(\mu_0,1/2)-sa_0}^{s}.
\]
Taking roots proves \eqref{eq:bonf}. Empty prime sets are immediate.

At an actual raw-minimal failure choose
\(s=\lfloor1/\eta\rfloor+1\) and \(z=M^{1/[2(s+1)]}\).
Chebyshev's estimate and \(M<16m^2\) give
\[
 E_s(k)=O_s\!\left(\frac{\sqrt M}{(\log M)^{s+1}}\right)=o_s(m).
\]
Therefore, in the sparse regime,
\begin{equation}
 \mu_{\rm low}:=\sum_{p\le z,\ w_p\ge\eta}\frac1p
 =O_{c,\eta}(\rho^{1/s}).\label{eq:lowharmonic}
\end{equation}
The rounding term in \eqref{eq:bonf} is essential. The argument does
not permit replacing \(z\) by the largest source prime and dropping
that term.

\section{Deleting exceptional sources by their weighted demand}

\begin{lemma}[Weighted monomial bound]
For a finite prime set \(Q\), let \(\mu_Q=\sum_{p\in Q}1/p\) and
\(\Omega_Q(n)=\sum_{p\in Q}v_p(n)\). For every integer \(k\ge1\),
\begin{equation}
 \sum_{\substack{n\le M\\\Omega_Q(n)\ge k}}X(n)
 \le M(k+\mu)h_k((1/p)_{p\in Q})
 \le M(k+\mu)\mu_Q^k.\label{eq:monomial}
\end{equation}
Here \(h_k\) is the complete homogeneous symmetric polynomial;
prime powers are counted with their repetitions.
\end{lemma}
\begin{proof}
Every integer on the left has a \(Q\)-smooth divisor \(d\) of total
degree \(k\). For each such distinct monomial,
\[
 \sum_{r\le M/d}X(dr)
 =\lfloor M/d\rfloor X(d)+\sum_{r\le M/d}X(r)
 \le(M/d)(k+\mu).
\]
The nonnegative union bound and summation over \(d\) prove the first
inequality. The expansion of \((\sum_Q1/p)^k\) contains every such
monomial with a coefficient at least one, proving the second.
\end{proof}

Apply this to the low substantial primes in \eqref{eq:lowharmonic},
with \(k=s+1\). Using \eqref{eq:sparsemean}, the weighted demand of all
sources with at least \(s+1\) such occurrences is
\[
 O_{c,\eta}(M\rho^{(s+1)/s})
 =O_{c,\eta}(m\rho^{1/s})=o(m).
\]
Every remaining source has at most \(s\) low substantial occurrences.
It has at most \(2s+1\) occurrences of primes exceeding
\(z=M^{1/[2(s+1)]}\), because \(2s+2\) such occurrences would have
product greater than \(M\). Thus it has at most
\begin{equation}
 K(\eta)=3s+1\le4s\label{eq:boundedocc}
\end{equation}
substantial occurrences in total. This bounds occurrences per source,
not the number of prime values in their union.

If the exceptional sources have weighted demand \(W_B=o(m)\), the
same original dual satisfies \(\Phi_{A\setminus B}=\Phi_A-W_B\).
For any subset \(U\) of its full hot set, original tight-row balance
also gives
\[
 \sum_pw_p\pos{V_p(U)-D_p(A\setminus B)}\le W_B.
\]
Hence both the objective and the positive row-excess normalization
survive this deletion uniformly. No new minimality or optimality of
the retained set is asserted.

\section{A local hot-divisor lemma and its exact scope}

Let \(a\) have load \(s_a=X(a)\), let \(u=\max_{q\mid a}w_q\), and
fix an occurrence of \(p\mid a\) with \(v=w_p>0\).
For any integer \(k\ge2\), if
\begin{equation}
 s_a>k+(k-1)(u-v),\label{eq:partitioncondition}
\end{equation}
there is a divisor \(d\mid a\) such that
\begin{equation}
 p\mid d,\quad X(d)>1,\quad d\le p^{1-1/k}a^{1/k}.
 \label{eq:partitiond}
\end{equation}
To see this, remove one occurrence of \(p\), and greedily form the
first \(k-1\) groups of remaining occurrences, stopping each just
after its weight exceeds \(1-v\). Each group's weight is at most
\(1-v+u\). Condition \eqref{eq:partitioncondition} ensures that the
remaining last group also has weight greater than \(1-v\). Assign
zero-weight occurrences to any group. The product of the \(k\)
group-products is \(a/p\); one is at most \((a/p)^{1/k}\). Multiply
that group by \(p\) to obtain \eqref{eq:partitiond}.

At an actual optimum its entire bank is hot, and active-row balance
therefore gives
\begin{equation}
 D_p(A)\ge\left\lfloor(M/p)^{1-1/k}\right\rfloor.
 \label{eq:partitiondemand}
\end{equation}
For a maximum-weight occurrence, \eqref{eq:partitioncondition} only
requires \(s_a>k\). In particular, every maximum-weight prime in a
source of load greater than two obeys
\(D_p(A)\ge\lfloor\sqrt{M/p}\rfloor\).

If \(Q_k\) is a set of qualifying active primes, then
\[
 \sum_{p\in Q_k}\frac{\log p}{p^{1-1/k}}
 \le\frac{2m\log M}{M^{1-1/k}}.
\]
This follows by multiplying \eqref{eq:partitiondemand} by \(\log p\),
using \(\lfloor t\rfloor\ge t/2\) for \(t\ge1\), and summing the
actual demands. These inequalities need not be small enough to
contradict the minimal class throughout \(M<16m^2\). They do not
pay for all core coordinates with one hinge budget.

\section{A stronger LP: whole prime-power components}

This is a separate possible route. Its necessary conditions must not
be assigned to a raw-minimal instance without a proof of transfer.
Let
\(N_A(p^j)=|\{a\in A:p^j\mid a\}|\). Replace the valuation constraints
in \eqref{eq:primal} by
\[
 \sum_{b\in I:p^j\mid b}z_b\ge N_A(p^j)
 \quad(1\le j\le\max_{a\in A}v_p(a)).
\]
The resulting optimum \(\atLP\) is at least \(\LP\). Its dual has
nonnegative level prices and monotone profiles
\[
 U(n)=\sum_p u_p(v_p(n)),\qquad
 u_p(e)=\sum_{j\le e}\alpha_{p,j}.
\]
Every profile may be capped at one without decreasing the objective.
Indeed, write \(U=U'+T\), with \(u'_p=\min(u_p,1)\). Then
\((U-1)_+=(U'-1)_++T\), because any positive tail already has a
capped contribution of one. The nonnegative prime-power increments
of \(T\) give \(\sum_A T\le\sum_{n\le M}T\le\sum_I T\).

At this LP's own minimal failure of a budget \(c>2\), every source
has \(U(a)>c\), hence at least three distinct active primes. All
targets with \(U>1\) are saturated, and each active level row is
exactly tight. If \(Q_q=q^{v_q(a)}\), the full bank of \(a/Q_q\)
is hot, since removing a whole prime profile costs at most one.
For any different active prime \(p\mid a\) and one of its active
levels \(p^j\mid a\), this gives
\begin{equation}
 N_A(p^j)\ge N_I(a/Q_q)\ge Q_q.
 \label{eq:componentm}
\end{equation}
In particular \emph{every} complete prime-power component of every
source is at most \(m\), including inactive components.

There is a stronger exact union inequality. For fixed \(a,p\), all
intersections of two distinct banks \(a/Q_q\), \(q\ne p\), are exactly
the \(a\)-multiples. Put
\(t_a=\lfloor M/a\rfloor\) and
\(S_a=\sum_{q\mid a}(Q_q-1)\). In every complete period of length
\(a\), their union contains exactly
\(1+S_a-(Q_p-1)\) positions. Thus
\begin{equation}
 N_A(p^j)\ge t_a\{1+S_a-(Q_p-1)\}.
 \label{eq:componentunion}
\end{equation}
Sum over one active level for each of the \(r_a\ge3\) active primes
in \(a\). Each demand is at most \(m\), while
\(\sum_{p\ {\rm active}}(Q_p-1)\le S_a\). Therefore
\begin{equation}
 \boxed{\quad
 \sum_{q\mid a}(q^{v_q(a)}-1)
 \le\frac{r_a}{r_a-1}\left(\frac m{\lfloor M/a\rfloor}-1\right)
 \le\frac32\left(\frac m{\lfloor M/a\rfloor}-1\right).
 \quad}\label{eq:componentsum}
\end{equation}
This is an actual-arithmetic union count, not a sum that repeatedly
spends the same capacity within a row. It has not been shown to
contradict all possible minimal instances.

The number of this LP's rows is now at most
\[
 \sum_{j\ge1}\pi(m^{1/j})
 \le\pi(m)+\sqrt m\log_2m=o(m).
\]
A cubic ambient gate \(M\ge8m^3\) also holds for the stronger integer
count-cover problem. Its complete proof, including the assignment
of whole prime-power components, is in the accompanying atom-LP
note. The raw quadratic gate cannot simply be imported: splitting
a prime power across two targets may satisfy valuation demand while
failing the separate level-count demand. Most importantly, none of
these stronger-LP reductions proves the required upper bound on its
optimum, nor a pointwise \(J_{\rm weak}\) for a fixed original raw dual.

\section{A verified failure of a same-threshold shortcut}

Bounded substantial occurrences per source do not imply a uniform
quadratic-error comparison of the two \emph{threshold-one} hinges.
The accompanying certificate constructs genuine integer intervals
with \(k\) unit-weight primes in \([P,1.1P]\), \(M=p_{\max}^4\), and
\begin{equation}
 \sum_{a\le M}(X(a)-1)_+-\sum_{b\in I}(X(b)-1)_+
 \ge\binom{k-1}{4}-\binom k2-3k.
 \label{eq:barrier}
\end{equation}
Each prefix integer has at most four support-prime occurrences,
including powers. The interval is centered at a CRT solution
\[
 C\equiv(p^5-p)/2\pmod{p^5}
 \quad\text{for every support prime }p.
\]
The proof uses exact inclusion-exclusion for the capped support count
and then restores the raw powers. It is supplied in full with an
integer-arithmetic verifier, rather than inferred from a numerical
sampling experiment.

For the explicit set of twenty primes from \(1009\) through \(1123\),
\(M=1123^4=1590446354641\). The independently checked certificate gives
\[
 \text{capped threshold-one deficit}=4360,\qquad
 \text{raw threshold-one deficit}=4321.
\]
In a growing close-prime family the lower bound in \eqref{eq:barrier}
is \(\Omega(k^4)\). The smallest support prime tends to infinity and
\(\sum_p1/(p-1)\) can tend to zero. Therefore these extra properties
do not rescue an \(O(k^2)\) same-threshold error bound.

\textbf{This is not a counterexample to \(J_{\rm weak}\).} The
requested inequality uses a source threshold above two and a target
threshold one, together with minimality and active-row balance. The
certificate does not satisfy that full set of hypotheses. Its role
is only to exclude a proposed intermediate lemma that would otherwise
silently replace the required joint argument.

\subsection{Why bounded active occurrences do not control all carriers}

There is an even more elementary obstruction to two tempting carrier
counts. Fix three active primes \(p_1<p_2<p_3\), all with weight one,
put \(Q=p_1p_2p_3\), and take
\[
 A=\{Qu:1\le u\le L,\ (u,Q)=1\},\qquad M=QL.
\]
Every source has exactly three active occurrences. If its inactive
cofactor is discarded, removing \(p_3\) gives the same active carrier
\(d=p_1p_2\) for \(\asymp_Q L\) sources. This is not a prime fibre:
\(\pi(M/d)=\pi(p_3L)=o_Q(L)\).

If instead one keeps the inactive factor, the carriers are
\(d_u=p_1p_2u\). At the target
\[
 b=p_1p_2\mathop{\rm lcm}\{u\le L:(u,Q)=1\},
\]
all these carriers divide \(b\). The remaining cofactor \(b/d_u\)
contains no active prime. A discount that examines only extra active
primes therefore assigns its full positive value to every carrier,
although \((X(b)-1)_+=1\). Their positive multiplicity is unbounded.

This example is also \emph{not} an actual optimal obstruction: its
entire \(p_1p_2\)-bank is hot and supplies at least \(p_3L\) units of
\(p_1\), exceeding \(D_{p_1}(A)\le L\). It demonstrates exactly why
bounded active degree alone does not justify either carrier count,
and why the actual joint demand balances must enter the missing proof.

\section{The precise missing estimate}

For the original normalized data of Section 1 there is an exact
identity with a single target budget:
\begin{equation}
 \Phi_G(w)=|U|+\sum_pw_p\{D_p(G)-V_p(U)\}
       -\sum_{b\in I\setminus U}\pos{X_w(b)-1}.
 \label{eq:residual}
\end{equation}
Replacing the signed residual demands by their coordinatewise
positive parts adds exactly \(\Xi=o(m)\). The missing statement is
an upper bound on \eqref{eq:residual} by
\((2+2\delta)|G|+o_\delta(m)\). None of the new lemmas bounds the
remaining \emph{positive} residual demands by that amount.

In the sparse branch, the reductions now provide all of the following
without assuming a special source class: a weighted mean at most
\(1+o(1)\); a bounded number of substantial occurrences per retained
source at each fixed weight threshold; substantial primes tending to
infinity; a derived weight gap with small objective loss; and the
original exact active-row balance for the retained hot targets.
The global number of substantial primes remains unbounded.

The unresolved step is a uniform \emph{joint} comparison for the
diffuse load and these substantial occurrences. Applying a conditional
estimate independently to each heavy factor pattern charges the
same target repeatedly. A theorem for finitely many exceptional
coordinates globally does not apply to a bound on occurrences in
each individual source. The band-deletion lemma resolves a numerical
parameter circularity, but does not resolve this overlap issue.
Its controlled linear loss also cannot be magnified by unspecified
pattern-dependent constants.

The branch in which \(m/M\) stays bounded away from zero has not been
excluded either. A separate companion note records an approximation
of its small-weight core by a constant using first and second moments;
it leaves a shifted heavy-profile hinge estimate unproved.

Accordingly, the proved facts in this report are reductions and exact
auxiliary inequalities. Neither \(J_{\rm weak}\), the stronger formula
(91) of the earlier report, the unrestricted coefficient-two
asymptotic, nor the exact \(2n\) bound is claimed. No new finite value
of \(g(n)\) is computed, and no additional restricted covering class
is offered as fulfillment of the stated goal.

\begin{thebibliography}{9}
\bibitem{previous} Research working note, \emph{Erd\H{o}s 708: Global
Dual Normalization and the Remaining Joint Inequality}, 23 November
2026. Preceding report and its reproducibility bundle.
\bibitem{wang} Interval-divisor-cover project, result catalogue and
retained source proof of the quadratic ambient theorem:
\url{https://github.com/science-prince/erdos-708-interval-divisor-covers/blob/main/docs/RESULTS.md}.
The theorem is used here as an inherited, previously audited input.
\bibitem{chen} H. Chen, \emph{Explicit upper bounds for the
Erd\H{o}s--Sur\'anyi function \(g(n)\)}, source repository:
\url{https://github.com/chy4pro/erdos-708-explicit-upper-bounds}.
The prime-power-count hinge formulation is also discussed in the
retained manuscript's remaining-problems section.
\bibitem{bloom} T. F. Bloom, \emph{Erd\H{o}s Problem 708},
\url{https://www.erdosproblems.com/708}. Problem statement, including
the separate asymptotic and exact coefficient-two questions.
\end{thebibliography}
\end{document}
